% Lines 808-1142 from original attention_transformers_notes_ghost.tex
% Chapter 5: Multi-Head Attention: Learning Diverse Representations

%==============================================================================
\section{Multi-Head Attention: Learning Diverse Representations}
%==============================================================================

\subsection{Motivation: Multiple Representation Subspaces}

Single attention allows the model to focus on one type of relationship at a time. However, natural language involves multiple types of relationships simultaneously:
\begin{itemize}
    \item Syntactic relationships (subject-verb, modifier-noun)
    \item Semantic relationships (synonyms, antonyms)
    \item Positional relationships (adjacent words, phrase boundaries)
    \item Long-range dependencies (coreference, discourse)
\end{itemize}

Multi-head attention enables the model to jointly attend to information from different representation subspaces.

\subsection{Mathematical Formulation}

Instead of performing a single attention function, multi-head attention:
\begin{enumerate}
    \item Projects queries, keys, and values $h$ times with different projections
    \item Performs attention in parallel for each projection
    \item Concatenates the results and projects again
\end{enumerate}

\begin{equation}
\text{MultiHead}(\mathbf{Q}, \mathbf{K}, \mathbf{V}) = \text{Concat}(\text{head}_1, ..., \text{head}_h)\mathbf{W}^O
\end{equation}

where each head is computed as:
\begin{equation}
\text{head}_i = \text{Attention}(\mathbf{Q}\mathbf{W}_i^Q, \mathbf{K}\mathbf{W}_i^K, \mathbf{V}\mathbf{W}_i^V)
\end{equation}
\begin{notationbox}
\textbf{Notation:}
\begin{itemize}
    \item $h$: Number of attention heads
    \item $\mathbf{W}_i^Q \in \mathbb{R}^{d_{model} \times d_k}$: Query projection for head $i$
    \item $\mathbf{W}_i^K \in \mathbb{R}^{d_{model} \times d_k}$: Key projection for head $i$
    \item $\mathbf{W}_i^V \in \mathbb{R}^{d_{model} \times d_v}$: Value projection for head $i$
    \item $\mathbf{W}^O \in \mathbb{R}^{hd_v \times d_{model}}$: Output projection
    \item $d_{model}$: Model dimension
    \item $d_k = d_v = d_{model}/h$: Dimension per head
\end{itemize}
\end{notationbox}

\subsection{Visualization of Multi-Head Attention}

\begin{figure}[H]
\centering
\begin{tikzpicture}[scale=0.7, every node/.style={scale=0.8}]
    % Input
    \node[draw, fill=gray!20, minimum width=3cm, minimum height=1cm] (input) at (0,0) {Input\\$\mathbf{X} \in \mathbb{R}^{n \times d_{model}}$};

    % Projection to heads
    \foreach \i/\color in {1/red,2/blue,3/green} {
        \node[draw, fill=\color!20, minimum width=2cm, minimum height=0.8cm] (qkv\i) at (-3+\i*3,2) {$\mathbf{Q}_{\i}, \mathbf{K}_{\i}, \mathbf{V}_{\i}$};
        \draw[arrow] (input) -- (qkv\i);

        % Attention
        \node[draw, fill=\color!40, minimum width=2cm, minimum height=0.8cm] (att\i) at (-3+\i*3,4) {Attention\\Head \i};
        \draw[arrow] (qkv\i) -- (att\i);

        % Output
        \node[draw, fill=\color!20, minimum width=1.5cm, minimum height=0.6cm] (out\i) at (-3+\i*3,6) {$\mathbf{Z}_{\i}$};
        \draw[arrow] (att\i) -- (out\i);
    }

    % Concatenation
    \node[draw, fill=yellow!20, minimum width=4cm, minimum height=0.8cm] (concat) at (3,7.5) {Concatenate};
    \foreach \i in {1,2,3} {
        \draw[arrow] (out\i) -- (concat);
    }

    % Final projection
    \node[draw, fill=orange!20, minimum width=3cm, minimum height=0.8cm] (proj) at (3,9) {Project $\mathbf{W}^O$};
    \draw[arrow] (concat) -- (proj);

    % Output
    \node[draw, fill=gray!20, minimum width=3cm, minimum height=1cm] (output) at (3,10.5) {Output\\$\mathbf{Z} \in \mathbb{R}^{n \times d_{model}}$};
    \draw[arrow] (proj) -- (output);

    % Title
    \node[above=0.5cm of output, font=\large\bfseries] {Multi-Head Attention Architecture};
\end{tikzpicture}
\caption{Multi-Head Attention with 3 Heads}
\label{fig:multihead-attention}
\end{figure}

\subsection{Implementation in PyTorch}

Let's implement multi-head attention following our strict coding guidelines:

\begin{lstlisting}[language=Python, caption=Multi-Head Attention Implementation]
from __future__ import annotations
from dataclasses import dataclass
import torch
import torch.nn as nn
import torch.nn.functional as F
from typing import Optional
import math


@dataclass(frozen=True, slots=True)
class AttentionConfig:
    d_model: int
    n_heads: int
    dropout: float = 0.1

    def __post_init__(self) -> None:
        if self.d_model % self.n_heads != 0:
            raise ValueError(
                f"d_model ({self.d_model}) must be divisible by "
                f"n_heads ({self.n_heads})"
            )
        if not 0 <= self.dropout < 1:
            raise ValueError("Dropout must be in [0, 1)")

    @property
    def d_k(self) -> int:
        return self.d_model // self.n_heads

    @property
    def d_v(self) -> int:
        return self.d_model // self.n_heads


class MultiHeadAttention(nn.Module):
    def __init__(self, config: AttentionConfig) -> None:
        super().__init__()
        self._config: AttentionConfig = config
        self._d_model: int = config.d_model
        self._n_heads: int = config.n_heads
        self._d_k: int = config.d_k
        self._d_v: int = config.d_v

        self._w_q: nn.Linear = nn.Linear(self._d_model, self._d_model)
        self._w_k: nn.Linear = nn.Linear(self._d_model, self._d_model)
        self._w_v: nn.Linear = nn.Linear(self._d_model, self._d_model)
        self._w_o: nn.Linear = nn.Linear(self._d_model, self._d_model)

        self._dropout: nn.Dropout = nn.Dropout(config.dropout)
        self._scale: float = math.sqrt(self._d_k)

        self._initialize_weights()

    def _initialize_weights(self) -> None:
        for module in [self._w_q, self._w_k, self._w_v, self._w_o]:
            nn.init.xavier_uniform_(module.weight)
            if module.bias is not None:
                nn.init.constant_(module.bias, 0.0)

    def forward(
        self,
        query: torch.Tensor,
        key: torch.Tensor,
        value: torch.Tensor,
        mask: Optional[torch.Tensor] = None
    ) -> torch.Tensor:
        batch_size: int = query.size(0)
        seq_len: int = query.size(1)

        Q: torch.Tensor = self._w_q(query)
        K: torch.Tensor = self._w_k(key)
        V: torch.Tensor = self._w_v(value)

        Q = self._reshape_for_attention(Q, batch_size)
        K = self._reshape_for_attention(K, batch_size)
        V = self._reshape_for_attention(V, batch_size)

        attention_output: torch.Tensor = self._scaled_dot_product_attention(
            Q, K, V, mask
        )

        attention_output = self._reshape_from_attention(
            attention_output, batch_size, seq_len
        )

        output: torch.Tensor = self._w_o(attention_output)
        return output

    def _reshape_for_attention(
        self,
        tensor: torch.Tensor,
        batch_size: int
    ) -> torch.Tensor:
        return tensor.view(
            batch_size, -1, self._n_heads, self._d_k
        ).transpose(1, 2)

    def _reshape_from_attention(
        self,
        tensor: torch.Tensor,
        batch_size: int,
        seq_len: int
    ) -> torch.Tensor:
        return tensor.transpose(1, 2).contiguous().view(
            batch_size, seq_len, self._d_model
        )

    def _scaled_dot_product_attention(
        self,
        Q: torch.Tensor,
        K: torch.Tensor,
        V: torch.Tensor,
        mask: Optional[torch.Tensor] = None
    ) -> torch.Tensor:
        scores: torch.Tensor = torch.matmul(Q, K.transpose(-2, -1)) / self._scale

        if mask is not None:
            scores = scores.masked_fill(mask == 0, -1e9)

        attention_weights: torch.Tensor = F.softmax(scores, dim=-1)
        attention_weights = self._dropout(attention_weights)

        context: torch.Tensor = torch.matmul(attention_weights, V)
        return context
\end{lstlisting}

\subsection{Theoretical Analysis of Multi-Head Attention}

\subsubsection{Representation Power of Multi-Head vs Single-Head}

Multi-head attention provably increases the representational capacity compared to single-head attention.

\textbf{Theorem (Representation Separation):} There exist functions computable by multi-head attention that cannot be efficiently computed by single-head attention with the same parameter budget.

\textbf{Proof:} Consider a sequence where we need to simultaneously:
\begin{enumerate}
    \item Track syntactic dependencies (e.g., subject-verb agreement)
    \item Model semantic similarity (e.g., word meaning relationships)
\end{enumerate}

A single attention head must choose a single similarity metric (projection), while multi-head can dedicate different heads to different similarity notions. Formally, let $f_1, f_2$ be two orthogonal similarity functions. Multi-head can compute:
\begin{equation}
\text{MH}(\mathbf{X}) = \text{Concat}(\text{Attn}(\mathbf{X}, f_1), \text{Attn}(\mathbf{X}, f_2))\mathbf{W}^O
\end{equation}

while single-head is restricted to a linear combination:
\begin{equation}
\text{SH}(\mathbf{X}) = \text{Attn}(\mathbf{X}, \alpha f_1 + \beta f_2)
\end{equation}

The non-linearity of attention makes these fundamentally different.

\subsubsection{Low-Rank Bottleneck Analysis}

A critical theoretical insight is that each head operates in a lower-dimensional subspace.

\textbf{Theorem (Rank Constraint):} For multi-head attention with $h$ heads and model dimension $d$, each head operates on a subspace of dimension $d/h$, creating a rank bottleneck.

\begin{notationbox}
\textbf{Notation:}
\begin{itemize}
    \item $h$: Number of attention heads
    \item $d$: Model dimension
    \item $d_k = d_v = d/h$: Dimension per head
    \item $\text{rank}(\cdot)$: Matrix rank
\end{itemize}
\end{notationbox}

\textbf{Analysis:} For each head $i$, the attention output is:
\begin{equation}
\text{head}_i = \text{Attention}(\mathbf{X}\mathbf{W}_i^Q, \mathbf{X}\mathbf{W}_i^K, \mathbf{X}\mathbf{W}_i^V)
\end{equation}

where $\mathbf{W}_i^Q, \mathbf{W}_i^K, \mathbf{W}_i^V \in \mathbb{R}^{d \times d_k}$. Thus:
\begin{equation}
\text{rank}(\text{head}_i) \leq \min(n, d_k) = \min(n, d/h)
\end{equation}

This creates a fundamental trade-off: more heads provide diversity but reduce per-head capacity.

\subsubsection{Optimal Number of Heads: Theoretical Analysis}

We derive the optimal number of heads from an information-theoretic perspective.

\textbf{Theorem (Optimal Head Count):} Under certain assumptions, the optimal number of heads $h^*$ satisfies:
\begin{equation}
h^* = \arg\max_h \left\{ h \cdot I(\mathbf{X}; \mathbf{Z}_{\text{head}}) - \lambda \cdot \frac{h}{d} \right\}
\end{equation}

where $I(\mathbf{X}; \mathbf{Z}_{\text{head}})$ is the mutual information between input and per-head output, and $\lambda$ is a capacity penalty.

\textbf{Intuition:} This balances two forces:
\begin{itemize}
    \item More heads → more diverse representations (first term)
    \item More heads → less capacity per head (second term)
\end{itemize}

\subsubsection{Attention Head Redundancy and Pruning Theory}

Not all attention heads contribute equally; some are redundant.

\textbf{Definition (Head Importance Score):} For head $i$:
\begin{equation}
\text{Importance}_i = \mathbb{E}_{\mathbf{x} \sim \mathcal{D}} \left[ \left\| \frac{\partial \mathcal{L}}{\partial \text{head}_i} \right\|_F \right]
\end{equation}

\textbf{Theorem (Redundancy Bound):} In a trained transformer with $h$ heads, at least $\lfloor h/4 \rfloor$ heads can typically be pruned with less than 1\% accuracy loss.

This surprising result suggests significant redundancy in multi-head attention, leading to recent work on dynamic head pruning.

\subsubsection{Multi-Head Attention as Ensemble Learning}

Multi-head attention can be viewed as an ensemble method with theoretical guarantees.

\textbf{Theorem (Variance Reduction):} Let $\sigma^2$ be the variance of a single head's output. With $h$ independent heads, the ensemble variance is reduced to:
\begin{equation}
\text{Var}[\text{MultiHead}] \approx \frac{\sigma^2}{h} + \mathcal{O}\left(\frac{1}{h^2}\right)
\end{equation}

This variance reduction partially explains the stability and generalization benefits of multi-head attention.

\subsection{What Different Heads Learn}

Research has shown that different attention heads often specialize in different types of relationships:

\begin{table}[H]
\centering
\caption{Typical Attention Head Specializations}
\begin{tabular}{l l}
\toprule
\textbf{Head Type} & \textbf{Learned Pattern} \\
\midrule
Positional & Attends to adjacent or nearby tokens \\
Syntactic & Tracks syntactic dependencies (e.g., subject-verb) \\
Semantic & Groups semantically related words \\
Rare tokens & Focuses on infrequent or important words \\
Delimiter & Attends to punctuation and sentence boundaries \\
Global & Broadly attends to many positions \\
\bottomrule
\end{tabular}
\end{table}
